%==============================================================================
%  Underground Design and Construction -- shared preamble.
%
%  Follows the Springer SNmono setup used for
%  Dynamic_Testing_and_Engineering_Applications_of_Brittle_Geomaterials.
%  There are no statement boxes and no colour anywhere in the book: material
%  that would once have been framed now runs in the main text, with a bold
%  run-in lead-in where the point needed a label.
%==============================================================================

% --- Springer Nature template packages ---------------------------------------
\usepackage{type1cm}
\usepackage{makeidx}
\usepackage{graphicx}
\usepackage{multicol}
\usepackage[bottom]{footmisc}
\usepackage{newtxtext}
\usepackage[varvw]{newtxmath}
\usepackage{bm}

% --- Content packages --------------------------------------------------------
\usepackage{amsmath}
% newtxmath already defines \Bbbk; release it before amssymb reloads it.
\let\Bbbk\relax
\usepackage{amssymb}
\usepackage{array}
\usepackage{booktabs}
\usepackage{tabularx}
\usepackage{longtable}[=v4.13]
\usepackage{multirow}
\usepackage{makecell}
\usepackage{ragged2e}
\usepackage{float}
\usepackage{enumitem}
\usepackage{siunitx}
\sisetup{mode=math}
\usepackage{microtype}
\usepackage{csquotes}
\usepackage{xurl}
\usepackage{tikz}
\usetikzlibrary{arrows.meta,positioning,fit,backgrounds,calc}

\graphicspath{{figures/}}

% --- Data figures: which vector format is included ---------------------------
% tools/export_twin_figures.py writes every data figure three times from the
% same drawing: PDF, SVG, and a PNG that exists only to be glanced at. The two
% vector files are siblings, not conversions -- neither is made from the other,
% and neither is made from a raster -- so the geometry is identical in both and
% the type is real type in both.
%
% PDF is the default route because pdflatex embeds it natively and as vector.
% The SVG route needs the svg package to call Inkscape at compile time, which
% means Inkscape on PATH and the whole document compiled with -shell-escape;
% Inkscape then writes a PDF and includes that. So the switch chooses which
% program converts the drawing, not whether the figure is vector: it is vector
% either way. Turn it on when the SVG is the file being edited by hand and the
% edits must reach the page without a re-export.
%
%   latexmk -pdf -shell-escape -outdir=build main.tex
%
% \datafigure[<width>]{<chapter>/<name>} -- no extension, both routes add it.
\newif\ifsvgfigures
\svgfiguresfalse

\ifsvgfigures
  \usepackage{svg}
  \svgpath{{figures/}}
  \newcommand{\datafigure}[2][\textwidth]{\includesvg[width=#1]{#2}}
\else
  \newcommand{\datafigure}[2][\textwidth]{\includegraphics[width=#1]{#2}}
\fi

% \renderfigure[<width>]{<scene>}{<labels>} -- a 3D render from figures/render/
% with labels typeset over it, so no text is burnt into the raster and the
% labels survive when a final render replaces the preview. Coordinates are
% fractions of the image: (0,0) bottom left, (1,1) top right.
%   \rlab[<anchor>]{<x,y of the point>}{<x,y of the text>}{<text>}  leader + label
%   \rtxt[<anchor>]{<x,y>}{<text>}                                  label only
\newcommand{\renderfigure}[3][0.9\textwidth]{%
  \begin{tikzpicture}[font=\footnotesize]
    \node[anchor=south west,inner sep=0] (img) at (0,0) {\includegraphics[width=#1]{render/#2}};
    \begin{scope}[x={(img.south east)},y={(img.north west)}]
      #3
    \end{scope}
  \end{tikzpicture}}
\newcommand{\rlab}[4][west]{%
  \draw[line width=0.3pt,black!80] (#2) -- (#3);
  \fill[black!80] (#2) circle[radius=0.6pt];
  \node[anchor=#1,inner sep=1.2pt,fill=white,fill opacity=0.85,text opacity=1,align=left] at (#3) {#4};}
\newcommand{\rtxt}[3][center]{%
  \node[anchor=#1,inner sep=1.2pt,fill=white,fill opacity=0.85,text opacity=1,align=center] at (#2) {#3};}

% --- Table conventions -------------------------------------------------------
\newcolumntype{Y}{>{\RaggedRight\arraybackslash}X}
\renewcommand{\arraystretch}{1.15}

% --- Figure palette ----------------------------------------------------------
% Shared with the Hydropower slide set and the Tri-HB Springer book. Colour is
% used for figures only; the running text and the statement environments remain
% unshaded, as elsewhere in this book.
\usepackage{xcolor}
%==============================================================================
%  SHARED FIGURE PALETTE
%  Identical, token for token, to the palette used by the Hydropower slide set
%  (RSE3241/Springer_Handbook/slide_palette_tikz.tex) and the Tri-HB Springer
%  book (Tri-HB/Springer_Book/figure_sources/palette.tex), so that a figure
%  carries between projects without recolouring. The values are not to be
%  adjusted here: change them in every project or in none.
%
%  The palette is built for a book read in colour on screen and printed in
%  black and white. Greyscale keeps lightness and discards hue, so lightness
%  carries the meaning and hue is confirmation. Strokes occupy a dark band and
%  fills a light one, and the two never meet: a line therefore always reads
%  dark against any fill, in colour, in greyscale, and under any colour-vision
%  deficiency.
%
%  Measured, relative to white:
%    strokes  CFAIL 8.0%  CMAIN 16.0%  CGOOD 23.9%  CLOSS 28.1%  luminance
%    fills    CWATER 62%  CPASS 72%    CSTOP 80%    CBAND 88%
%    contrast CFAIL 8.08:1  CMAIN 5.01:1  CGOOD 3.64:1  CLOSS 3.17:1 on white
%
%  THREE RULES
%   1. Colour is never the only carrier of meaning. Adjacent members of each
%      band are only about 4 L* apart -- CGOOD against CLOSS, and the three
%      light fills against one another, are close to indistinguishable in a
%      greyscale print. Pair colour with dash pattern, numbering or position
%      so the figure survives a photocopy.
%   2. CWATER is a fill, never a line: at 1.57:1 on white it fails the 3:1
%      graphics-contrast floor, and it is the token that collapses into green
%      under tritanopia.
%   3. In TikZ a bare colour name after fill= overrides the fill. Write
%      [fill=CBAND, draw=CLOSS], never [fill=CBAND, CLOSS].
%
%  Check every new figure in greyscale before committing it:
%      pdftocairo -png -gray -r 150 figure.pdf out
%==============================================================================

% ---- strokes: lines, arrows, labels. All >= 3:1 on white. -------------------
\definecolor{CFAIL}{HTML}{A50000}   %  8.0% grey, 8.08:1 -- failure, critical
\definecolor{CMAIN}{HTML}{007D76}   % 16.0% grey, 5.01:1 -- principal series
\definecolor{CGOOD}{HTML}{009C1C}   % 23.9% grey, 3.64:1 -- admissible, pass
\definecolor{CLOSS}{HTML}{C18500}   % 28.1% grey, 3.17:1 -- threshold, limit

% ---- fills: areas, bands, status boxes -------------------------------------
\definecolor{CWATER}{HTML}{87DAF8}  % 62.1% grey -- cold / water / slow
\definecolor{CPASS}{HTML}{B5E8CC}   % 71.9% grey -- admissible region
\definecolor{CSTOP}{HTML}{F8E2E2}   % 79.8% grey -- excluded region
\definecolor{CBAND}{HTML}{FFF1C5}   % 88.2% grey -- caution band

% ---- neutrals: structure that carries no category ---------------------------
\definecolor{CINK}{gray}{0.22}      % axes, frames, primary label text
\definecolor{CTEXT}{gray}{0.34}     % secondary annotation
\definecolor{CRULE}{gray}{0.78}     % grid lines, light rules
\definecolor{CBOX}{gray}{0.95}      % neutral panel fill

% ---- rate ramp: ordered cold-to-warm scale, built from the tokens above -----
\colorlet{rateslow}{CMAIN}
\colorlet{rateslowlt}{CWATER}
\colorlet{ratefastlt}{CSTOP}
\colorlet{ratefast}{CFAIL}

% ---- cross-project compatibility -------------------------------------------
%  Aliases used by figures written for the slide set and the Tri-HB book, so
%  that source copied from either compiles here unchanged. New figures in this
%  book use the tokens above directly.
\colorlet{CDARK}{black}
\colorlet{CGONE}{black!80}
\colorlet{CGTWO}{black!50}
\colorlet{CGTHREE}{black!12}
\colorlet{CACC}{CMAIN!80!black}
\colorlet{slate}{CINK}
\colorlet{steel}{CMAIN}
\colorlet{clay}{CFAIL}
\colorlet{ochre}{CLOSS}
\colorlet{textgrey}{CTEXT}
\colorlet{rulegrey}{CRULE}
\colorlet{boxgrey}{CBOX}
\colorlet{zonelite}{CBOX}
\colorlet{zonemid}{CRULE}
\colorlet{zonedark}{CTEXT}
\colorlet{teal}{CGOOD}


% --- TikZ workflow diagram styles --------------------------------------------
% Size and role are orthogonal: a node takes one size style and one role style,
% for example [smallworkflowbox, wfstep].
%
% The house palette carries two hues, and these diagrams carry four roles, so
% role is encoded as STROKE HUE times FILL rather than by hue alone:
%
%   wfinput   grey stroke,      white fill   what the sequence starts from
%   wfstep    malachite stroke, white fill   a computation
%   wfadopt   malachite stroke, tint fill    the value that enters the design
%   wfcheck   carmine stroke,   tint fill    a comparison against a threshold
%   feedbackarrow  carmine dashed            the path back when the check fails
%
% Malachite and carmine are close in lightness, so in a greyscale print the
% hues read alike and the FILL carries the distinction: an empty box is a step,
% a tinted box is an outcome. Steps are also numbered and the return path is
% dashed, so no diagram depends on colour alone. See palette.tex, rule 1.
%
% Figure text is set in Computer Modern per the house style. Mathematics inside
% a TikZ node still takes the document's newtxmath, because the font is fixed
% by the Springer class; that mixture is transitional and resolves as figures
% are redrawn outside TikZ, where CM mathematics is native.
\tikzset{
  wffont/.style={font=\fontfamily{cmr}\selectfont},
  workflowbox/.style={
    rectangle, rounded corners, draw=QGREY, fill=white, wffont,
    thick, align=center, text width=2.8cm, minimum height=1.65cm, inner sep=4pt},
  smallworkflowbox/.style={
    rectangle, rounded corners, draw=QGREY, fill=white, wffont,
    thick, align=center, text width=2.65cm, minimum height=1.45cm, inner sep=3pt},
  tinyworkflowbox/.style={
    rectangle, rounded corners, draw=QGREY, fill=white, wffont,
    thick, align=center, text width=2.45cm, minimum height=1.45cm, inner sep=3pt},
  wfinput/.style={draw=QGREY, fill=white, text=QINK},
  wfstep/.style={ draw=QMAL,  fill=white, text=QINK},
  wfadopt/.style={draw=QMAL,  fill=QMALT, text=QINK, very thick},
  wfcheck/.style={draw=QCAR,  fill=QCART, text=QINK},
  workflowarrow/.style={-{Latex[length=2mm]}, thick, draw=QINK},
  feedbackarrow/.style={-{Latex[length=2mm]}, thick, dashed, draw=QCAR}
}

% --- Draft aids --------------------------------------------------------------
% Set \draftmodefalse before submission to hide all editorial markers.
\newif\ifdraftmode
\draftmodetrue

\newcommand{\todo}[1]{\ifdraftmode\textbf{[TODO: #1]}\fi}
\newcommand{\source}[1]{\ifdraftmode\marginpar{\scriptsize #1}\fi}

% Outline-stage scope note: states what a section will cover and which method
% or source governs it. Renders as plain italic; no colour.
\newcommand{\scopenote}[1]{\ifdraftmode
  \par\smallskip\noindent{\small\itshape #1}\par\smallskip\fi}

% --- Statement boxes ---------------------------------------------------------
% None. The material that was formerly boxed now runs in the main text, with a
% bold run-in lead-in where the point needed a label. Springer's own box
% environments remain available through springer_book.sty if ever required.

% --- Notation ----------------------------------------------------------------
\newcommand{\sci}{\sigma_{ci}}                % UCS of intact rock
\newcommand{\scm}{\sigma_{cm}}                % rock mass strength
\newcommand{\sone}{\sigma_1}
\newcommand{\stwo}{\sigma_2}
\newcommand{\sthree}{\sigma_3}
\newcommand{\sv}{\sigma_v}                    % vertical in-situ stress
\newcommand{\sh}{\sigma_h}                    % horizontal in-situ stress
\newcommand{\pzero}{p_0}                      % far-field pressure
\newcommand{\pin}{p_i}                        % internal support pressure
\newcommand{\mb}{m_b}
\newcommand{\mi}{m_i}
\newcommand{\GSI}{\mathrm{GSI}}
\newcommand{\RMR}{\mathrm{RMR}}
\newcommand{\RQD}{\mathrm{RQD}}
\newcommand{\Erm}{E_{m}}                      % rock mass modulus
\newcommand{\Ei}{E_i}                         % intact rock modulus
\newcommand{\Dfac}{D}                         % disturbance factor

\newcommand{\etal}{et al.}
\newcommand{\code}[1]{\texttt{#1}}

% --- Companion open-access application ---------------------------------------
% The book is tool-neutral: no chapter depends on a proprietary program. Where a
% calculation is more convenient run than done by hand, the text points to the
% open-access companion instead. Both the name and the address are set here, so
% that moving or renaming the repository is a one-line change.
\newcommand{\companionname}{the companion open-access application}
\newcommand{\companionurl}{https://github.com/qbzhang666/RSE3010}
\newcommand{\companionlink}{\expandafter\url\expandafter{\companionurl}}
\newcommand{\companion}{\companionname\ (Appendix~\ref{app:software})}
